\documentclass[12pt,a4paper]{article}
% Preamble bersama untuk kedua laporan TK 2 (soal1 dan soal2).
% Mengikuti "Petunjuk Penulisan Dokumen Project Anum":
%   - A4, font standar LaTeX (Latin Modern), 12pt, margin >= 2.54 cm
%   - judul, rangkuman, daftar isi TIDAK dinomori; bagian lain dinomori dari 1
%   - tabel/gambar/kode/persamaan dinomori mengikuti bagian: Tabel 3.1, Gambar 3.2, Kode 3.1
\usepackage[T1]{fontenc}
\usepackage[utf8]{inputenc}
\usepackage[indonesian]{babel}
\usepackage{lmodern}
\usepackage[a4paper,margin=2.54cm]{geometry}
\usepackage{amsmath,amssymb,booktabs,array,graphicx,float,microtype}
\usepackage{algorithm,algpseudocode}
\usepackage{listings}
\usepackage[hidelinks]{hyperref}
\emergencystretch=1em

\numberwithin{equation}{section}
\counterwithin{table}{section}
\counterwithin{figure}{section}
\counterwithin{algorithm}{section}
\floatname{algorithm}{Kode}
\renewcommand{\lstlistingname}{Kode}
\algrenewcommand\algorithmicrequire{\textbf{Masukan:}}
\algrenewcommand\algorithmicensure{\textbf{Keluaran:}}
\setlength{\belowcaptionskip}{6pt}
\lstset{language=Python, basicstyle=\ttfamily\footnotesize, breaklines=true,
  columns=fullflexible, keepspaces=true, showstringspaces=false, frame=single}

\newcommand{\R}{\mathbb{R}}
\newcommand{\T}{\mathsf{T}}
\newcommand{\norm}[1]{\left\lVert #1\right\rVert_2}
\newcommand{\bx}{\mathbf{x}}
\newcommand{\bd}{\mathbf{d}}
\newcommand{\bg}{\mathbf{g}}
\newcommand{\bs}{\mathbf{s}}
\newcommand{\by}{\mathbf{y}}

% \input biasa di dalam tabular gagal bila diikuti \bottomrule; pakai primitif \@@input
\makeatletter\newcommand{\isitabel}[1]{\@@input #1 }\makeatother

% Grafik dicari di report/figures/ dan results/ (relatif terhadap folder report/)
\graphicspath{{figures/}{../results/}}

\begin{document}

% ---- Judul (1 baris, tidak dinomori) ----
\begin{center}
{\Large\bfseries Soal 2: Prediksi Data Hilang dengan Interpolasi Polinomial dan Spline}\par\medskip
Kelompok A9 -- TK 2 Analisis Numerik
\end{center}

% ---- Rangkuman (~1/2 halaman, tidak dinomori). Ditulis TERAKHIR. ----
\section*{Rangkuman}
% TODO: masalah, model (polinomial d=2..9, spline linear, spline kubik natural), galat, FLOP, waktu, rekomendasi.

\setcounter{tocdepth}{1}
\tableofcontents
\clearpage

\section{Pendahuluan}
% TODO (ditulis terakhir).

\section{Data}\label{sec:data}
% Tabel 2 (latih) dan Tabel 3 (uji) dari soal; pemetaan x = 2t/9 - 1.

\section{Interpolasi Polinomial Basis Standar}\label{sec:poly}
% Tugas (i): indeks i_j = floor(9j/d + 1/2), Vandermonde, Gauss pivot parsial, Horner; tabel koefisien semua derajat.

\section{Spline Linear}\label{sec:linspline}
% Tugas (ii): koefisien per segmen dalam koordinat lokal t - t_i.

\section{Spline Kubik Natural}\label{sec:cubspline}
% Tugas (iii): sebutkan algoritma (mis. Sauer); koefisien per segmen dalam t - t_i.

\section{Prediksi dan Galat}\label{sec:galat}
% Tugas (iv): tabel z_hat, z, |z_hat - z|, norma l2 untuk tiap model.

\section{Visualisasi}\label{sec:visual}
% Tugas (v): Gambar 1 (kurva poli + spline, 201 titik, penanda latih/uji), Gambar 2 (galat per titik uji; norma l2 vs derajat + 2 garis mendatar).

\section{Jumlah Operasi (FLOP)}\label{sec:flop}
% Tugas (vi): hitung pembentukan vs evaluasi lima titik uji; tunjukkan cara menghitung.

\section{Waktu Eksekusi}\label{sec:waktu}
% Tugas (vii): perf_counter; median 5 rangkaian x 100 panggilan; sebut lingkungan eksekusi + satuan.

\section{Analisis}\label{sec:analisis}
% Tugas (viii): tren galat vs derajat, tengah vs ujung, FLOP vs waktu, pilih satu model + alasan.

\section{Kesimpulan}
% TODO: 1--2 paragraf.

\begin{thebibliography}{9}
\renewcommand{\refname}{Referensi}
\bibitem{sauer} T. Sauer, \emph{Numerical Analysis}, edisi ke-3. Pearson, 2018, Bab 3.
% TODO: tambah referensi yang benar-benar dipakai (bukan slide kuliah).
\end{thebibliography}

\appendix
\section{Lampiran}
% Kode program (bukan pseudocode), tabel koefisien rinci.
\end{document}
